\scanpage{3}
132.

\uline{\S24. Explicit description of pseudo-rational double points.}

Let $R$ be a two-dimensional normal local ring of multiplicity two having a pseudo-rational singularity. Since $\lambda(m^2)=4$ (Corollary 23.3) every minimal basis of the maximal ideal $m$ of $R$ consists of three elements. We shall classify $R$ by studying its behaviour under successive quadratic transformations, and by relating this behaviour to certain conditions involving, more or less explicitly, a suitable basis $\{x,y,z\}$ of $m$. We also prove converse statements of the type: ``If $R$ is any local ring, with maximal ideal $m$ generated by elements $x,y,z$ satisfying $\ldots$ then $R$ is a two-dimensional normal local ring having a (pseudo-) rational singularity.'' In other \textbf{words} the conditions to be introduced \uline{characterize} (pseudo-) rational ``double points.''

Basically, the idea is to take a two-dimensional local ring of multiplicity two whose maximal ideal is generated by three elements, say $x,y,z$, to subject this ring to a succession of quadratic transformations, and to see what conditions on $x,y,z$, guarantee that the resulting rings are all \uline{normal} (cf. Proposition 23.5). This approach involves a detailed and rather tedious examination of numerous cases. For orientation, the reader may analyse a ring of the form
\[
k[[X,Y,Z]]/(Z^2-f(X,Y))
\]
($k$-a field) from this point of view.
\[
*\qquad\qquad *\qquad\qquad *
\]
\mspage{2}
